\subsection{\texorpdfstring{\mymethodmlp}{A3-MLP}}
\label{sec:method:ffn}

The non-linear activation function in \mlp{} component prohibits us from directly applying SVD.
Instead, we first derive an objective for minimizing the MLP output error,
and uses CUR decomposition~\citep{CUR-matrix-decompositions-for-improved-data-analysis} to find the low-rank form of MLP weights.

% \vspace{-0.5em}
\begin{problem}[Minimization of \mlp{} output error]\label{problem:mlp}
Given a pretrained down
% \gc{what does down mean here?} 
projection layer $\mYffn = \mXdown \Wdown$ in MLP
and its approximated low-rank form $\mYffnapprox = \mXdownapprox \Wdownapprox = \mXdown \mU \Wdown$, 
minimizing the MLP output error is to minimize the following loss: 
% \begin{equation}
%     \mathrm{arg min}_{\mU=\mathrm{diag}(u_1, u_2, \dots, u_\dinter)} \|\mXdown \mU \Wdown - \mXdown\Wdown\|_F^2
%     \quad \text{s.t.} \quad \text{rank}(\mU)=r\;,
% \end{equation}
\begin{equation}
\begin{aligned}
    \underset{\mU=\operatorname{diag}(u_1,\dots,u_\dinter)}{\mathrm{arg\,min}}
    \;& \|\mXdown \mU \Wdown - \mXdown \Wdown\|_F^2 \\
    \text{s.t.} \quad
    & \operatorname{rank}(\mU) = r \; .
\end{aligned}
\end{equation}

where $\mXdown \in \sR^{\ldown \times \dinter}$ is the matrix of intermediate activation vectors,
and $\mU \in \sR^{\dinter \times \dinter}$ is a diagonal matrix determining which $r$ columns of $\Wdown$ to keep.

% where $\vxdown \in \sR^{\dinter}$ is the intermediate activation vector from the input space $\sX_{\text{d}} \subseteq \sR^{\dinter}$
% and $\mU \in \sR^{\dinter \times \dinter}$ is a diagonal matrix determining which $r$ columns of $\Wdown$ to keep.
\end{problem}

%We force $\mU$ to be a diagonal matrix,
%because reducing the intermediate size $\dinter$ to $r$ means
%reducing the rank of $\Wdown$ to $r$, and $\Wup$, $\Wgate$, and .
%such that the zero elements on the diagonal determine which columns of $\vxdown$ (the rows of $\Wdown$) to remove
%Note that $\mU$ also keep $r$ columns of $\vxdown$, $\Wgate$ and $\Wup$ accordingly.

%\vspace{0.5em}
\begin{lemma}[Equivalent form of Problem \ref{problem:mlp}]
    \label{lemma:a3:mlp:obj:derived}
    The objective in Problem~\ref{problem:mlp} is equivalent to the following error on the calibration dataset:
    % \begin{equation}
    %     \mathrm{argmin}_{\mU=\mathrm{diag}(u_1, u_2, \dots, u_\dinter)} \|\mR_{\sX_{\text{d}}\sX_{\text{d}}}^{\frac{1}{2}} \mU \Wdown - \mR_{\sX_{\text{d}}\sX_{\text{d}}}^{\frac{1}{2}} \Wdown\|_F^2 
    %     \quad \text{s.t.} \quad \text{rank}(\mU)=r\;,
    % \end{equation}
    \begin{equation}
    \begin{aligned}
    \underset{\mU = \operatorname{diag}(u_1, u_2, \dots, u_{\dinter})}{\arg\min}
    \quad
    & \left\|
    \mR_{\sX_{\text{d}}\sX_{\text{d}}}^{\frac{1}{2}}
    \mU \Wdown
    -
    \mR_{\sX_{\text{d}}\sX_{\text{d}}}^{\frac{1}{2}}
    \Wdown
    \right\|_F^2 \\
    \text{s.t.} \quad
    & \operatorname{rank}(\mU) = r \; .
    \end{aligned}
    \end{equation}
    

    % \begin{equation}
    % \begin{aligned}
    %     \underset{\mU=\mathrm{diag}(u_1,\dots,u_\dinter)}{\mathrm{arg\,min}}
    %     \;& \|\mXdown \mU \Wdown - \mXdown \Wdown\|_F^2 \\
    %     \text{s.t.}\;& \mathrm{rank}(\mU) = r \;.
    % \end{aligned}
    % \end{equation}
    where $\mR_{\sX_{\text{d}}\sX_{\text{d}}} := \frac{1}{\ldown}\mXdown^T \mXdown$ is the autocorrelation matrix of the calibration set $\mXdown$.
\end{lemma}
The derivation of Lemma~\ref{lemma:a3:mlp:obj:derived} is in~\Cref{sec:appendix:derivation:mlp:objective}. This CUR approximation is a well-studied NP-hard problem, and various CUR methods have been proposed~\citep{optimal-cur-matrix-decompositions, Fast-Monte-Carlo-Algorithms-for-Matrices-I:-Approximating-Matrix-Multiplication}.
We thus pick a simple but effective solution from~\cite{Fast-Monte-Carlo-Algorithms-for-Matrices-I:-Approximating-Matrix-Multiplication} and name this approach \mymethodmlp{}.

Following~\cite{Fast-Monte-Carlo-Algorithms-for-Matrices-I:-Approximating-Matrix-Multiplication}, we build $\mU$ by
